\section{Implications: Affective Practices and Interaction Qualities}
\label{sec:ch3-higher-order}

We now introduce \textit{Affective Practices} (AP1---AP6) and \textit{Interaction Qualities} (IQ1---IQ6), explaining how they were derived. We then share each practice and its enabling interaction qualities, grounded in participant examples.

\subsection{Deriving Affective Practices and Interaction Qualities}
Our five themes trace how affective experience unfolds, beginning with bodily perception and anticipation, to situated adjustment and interpretations of automation. To make these themes actionable for design, we introduce \textbf{Affective Practices (APs)}: recurring ways occupants \textit{sense, interpret, and regulate} their experience in smart environments, often socially and in relation to system behaviour.\footnote{For readability, we use the term \textbf{the system} to refer to the underlying smart building logic, including its sensing, thresholds, and policies that shape environmental behaviour.}

We derived them from our themes in Study 2 (W1---W5 in \autoref{subsec:s2-analysis}) by identifying recurrent patterns of affective experience. Each practice emerges directly from one or more themes: W1’s embodied sensing supports AP1 (\textit{Orient and Appraise}) and AP2 (\textit{Select a Micro-Condition}); W2 and W3’s attention to social atmosphere and spatial purpose shape AP3 (\textit{Social Attunement and Coordination}); W4’s trajectories of drift, escalation, and pacing underpin AP4 (\textit{Modulate and (De)Escalate}) and AP6 (\textit{Settle and Stabilise}); and W5’s interpretive work around automation informs AP5 (\textit{Probe and Align with System Behaviour}). This mapping, summarised in \autoref{tab:ap-iq-summary}, keeps the practices empirically anchored while enabling system-facing design implications.

Analytically distinct, the practices are not mutually exclusive. Occupants often move between them, orienting, attuning socially, and adjusting in overlapping ways. Their value lies in foregrounding different loci of affective sense-making that smart environments can support, rather than enforcing rigid categorisation.

% what are interaction qualities
Supporting the APs, we also developed \textbf{Interaction Qualities (IQs)} using ``what if?'' scenarios: system-facing capacities that describe how buildings scaffold or respond to these practices. Together, APs and IQs offer a new set of goals and tools for designing novel interactions with(in) smart buildings.



\subsection{AP1---Orient and Appraise}
We define \textit{Orient and Appraise} as the initial perceptual process through which occupants develop an embodied sense of whether a space aligns with their needs. This practice unfolds through environmental scanning---attending to spatial and social cues---as well as tuning into internal bodily states such as posture, visual strain, or thermal comfort. It produces a ``soft judgment'': \textit{Is this good enough, for now, for what I want to do?} This orienting typically precedes any active interaction with the space. Participants noticed qualities like light direction, seating angles, or degrees of exposure. In response, they formed a somatic sense of fit that shaped their decision to settle, tweak posture, or move on. Unlike evaluative models of comfort, this is a pre-reflective, felt alignment that might suffice briefly.

\paragraph{\textbf{Enabling Interaction Quality---Somatic Legibility}}
To support this practice, smart building technologies must embody \textbf{Somatic Legibility}---the system’s capacity to help occupants perceive, interpret, and act on their bodily responses to the environment, without requiring explicit commands. First, systems can surface how others have responded to a space. For instance, a hallway lounge might project residual silhouettes of recent sitting postures using thermal data, evoking whether people tended to slouch, tense up, or relax. Second, systems can surface somatic cues with appropriate subtlety. When ambient CO\textsubscript{2} or thermal readings cross thresholds, a space might momentarily shift its colour temperature, not as an alarm, but as a somatic prompt to notice whether the space still ``feels right''. 

% \begin{tcolorbox}[colback=blue!5!white, colframe=blue!60!black, title=Enabling Interaction Quality: Somatic Legibility]
% Somatic Legibility is the system’s capacity to help occupants perceive, interpret, and act on their own bodily responses to the environment, without requiring explicit commands or constant attention.
% \end{tcolorbox}

\designbox
{Enabling Interaction Quality: Somatic Legibility}
{Somatic Legibility is the system’s capacity to help occupants perceive, interpret, and act on their own bodily responses to the environment, without requiring explicit commands or constant attention.}
{Boxed text presenting a conceptual definition}


\subsection{AP2---Select a Micro-Condition}
\textit{Select a Micro-Condition} is the practice of choosing a spatial niche---a bounded configuration shaped by enclosure, light, acoustic tone, airflow, or social proximity that aligns with one’s current needs. This is not a search for perfection, but a pragmatic act of tuning. Occupants often settle for ``good enough for now,'' balancing bodily state, task demands, and environmental trade-offs. Importantly, people rarely seek ``a room''; they seek a space that feels right.

\paragraph{\textbf{Enabling Interaction Quality---Spatial Niche Discovery}}
\textit{Spatial Niche Discovery} is the system’s capacity to surface contrasts, previews, and transitions between adjacent spaces, helping occupants sense suitable spaces. First, interaction should raise awareness of nearby options. Edge lighting embedded along shared desks might change hue based on real-time soundscapes, highlighting quieter zones at the periphery. A hallway panel might show live thermal-acoustic gradients across a floor, guiding users toward calmer spaces. Unlike fixed signage or occupancy sensors, these cues expose felt differences between niches, not just availability.

Second, interaction should enable interpretation each niche's affordances. As occupants move through the space, ambient tags or overlays on walls, furniture, or mobile devices could suggest temporal suitability (e.g. ``low distraction for 15 min''). These framings help occupants weigh trade-offs without extensive configuration or environmental literacy. Finally, systems should support transitions between niches. As someone moves between spots, the environment might soften background noise in the new area using directional sound masking. Such handovers can help occupants re-settle, treating spatial selection as a fluid negotiation of fit.

% \begin{tcolorbox}[colback=blue!5!white, colframe=blue!60!black, title=Enabling Interaction Quality: Spatial Niche Discovery]
% Spatial Niche Discovery is the system’s capacity to surface contrasts, previews, and low-friction transitions between adjacent spatial niches—helping occupants sense which space might suit them in the moment.
% \end{tcolorbox}

\designbox
{Enabling Interaction Quality: Spatial Niche Discovery}
{Spatial Niche Discovery is the system’s capacity to surface contrasts, previews, and low-friction transitions between adjacent spatial niches—helping occupants sense which space might suit them in the moment.}
{Boxed text presenting a conceptual definition}


\subsection{AP3---Social Attunement and Coordination}
\textit{Social Attunement and Coordination} is the practice of interpreting and participating in the informal social fabric of a space. Occupants observe one another to learn what kinds of activities are appropriate where and align their behaviour with local expectations. These expectations often take the form of ``soft norms''---informal understandings that guide behaviour in shared environments without being explicitly stated. Such norms are also shaped collectively over time. The same room may feel like a quiet focus zone in the morning and a lively discussion zone in the afternoon. In ambiguous or mixed-use zones, this interpretive work becomes especially important for feeling legitimate and welcome.

\paragraph{\textbf{Enabling Interaction Quality---Surfacing Social Cues}}
Technologies must prioritise \textbf{Social Cue Surfacing}---helping occupants perceive tacit cues about the evolving social environment. This quality reflects the vernacular, participatory nature of how people come to understand space. First, systems should surface aggregate social signals such as ambient conversational density, movement, or group cohesion, without exposing identities or tracking individuals. For example, an overhead light gradient could reflect sound patterns in a shared atrium: growing warmer as group discussions cluster, or cooler as ambient activity quiets down.

Second, interaction should support interpretation. Rather than binary labels like “quiet zone,” systems might display detailed atmosphere tags---such as ``murmur-type ,'' or “collaborations in progress” that reflect current use without directing. Finally, systems should enable behavioural adjustment that aligns with social changes. For instance, wearable prompts or glance-able ambient icons might socially sparse niches (``lower activity near window alcove''). The goal is not to pre-empt social coordination, but to provide information that can help people participate in a shared sense of how a space is currently being lived.

% \begin{tcolorbox}[colback=blue!5!white, colframe=blue!60!black, title=Enabling Interaction Quality: Surfacing Social Cues]
% Social Cue Surfacing helps occupants perceive tacit cues about the evolving social environment and reflects the vernacular, participatory nature of how people come to understand space.
% \end{tcolorbox}

\designbox
{Enabling Interaction Quality: Surfacing Social Cues}
{Social Cue Surfacing helps occupants perceive tacit cues about the evolving social environment and reflects the vernacular, participatory nature of how people come to understand space.
}
{Boxed text presenting a conceptual definition}


\subsection{AP4---Modulate and (de)Escalate}
\textit{Modulate and (De)Escalate} is the regulation of spatial fit during ongoing use, once initial settlement has been achieved, but misalignment begins to surface. Unlike \textit{Orient and Appraise}, which addresses first impressions, this practice takes place over time. Occupants' fit (with a space) may drift gradually or break suddenly due to sudden environmental change (daylight being blocked by automated blinds). In such moments, occupants engage in paced, proportionate, and often reversible actions to preserve task continuity: modulating the environment, escalating only when necessary, and de-escalating when the stressor passes.

\paragraph{\textbf{Enabling Interaction Quality---(De)Escalation-Aware Interaction}}  
\textit{(De)Escalation-Aware Interaction} refers to the system’s capacity to offer \textit{watchful support} in the form of minimally suggestive, reversible, and situationally attuned cues that assist occupants during regulatory drift without overriding their intent.

In shared environments, this can be through \textbf{scene technology} embedded in everyday modalities. Personal headphones, already widely used in open-plan workspaces, can deploy sound-based scenes: muffled forest tones (\textit{Reset}), low-frequency wind textures (\textit{Soften}), or binaural beats (\textit{Focus}) that support affective recalibration. These scenes are scoped in time and intensity, layering over the occupant’s experience without altering the shared environment. Watchful support emerges when the system detects signs of escalation or de-escalation, such as repeated adjustments, posture shifts, or prolonged scanning. Rather than triggering global changes, the interface might suggest: \textit{``Try Focus scene in your headphones for 20 min?''}  or, if the occupant appears to be searching for comfort, the system could propose: \textit{``Softer spot nearby, take your scene with you.''}  

% \begin{tcolorbox}[colback=blue!5!white, colframe=blue!60!black, title=Enabling Interaction Quality: (De)Escalation-Aware Interaction]
% (De)Escalation-Aware Interaction offers watchful support in the form of minimally suggestive, reversible, and situationally attuned cues that assist occupants during regulatory drift without overriding their intent.
% \end{tcolorbox}

\designbox
{Enabling Interaction Quality: (De)Escalation-Aware Interaction}
{(De)Escalation-Aware Interaction offers watchful support in the form of minimally suggestive, reversible, and situationally attuned cues that assist occupants during regulatory drift without overriding their intent.}
{Boxed text presenting a conceptual definition}


\subsection{AP5---Probe and Align with Smart Building Technology}
We define \textit{Probe and Align} as the situated practice through which occupants interpret, test, and revise their understandings of how smart environments behave. Often, smartness unfolds in the background via sensors, thresholds, and policies that act without visible triggers or explanations~\cite{saputra2023smart, daissaoui2020iot}. Occupants make sense of these behaviours as they go: observing patterns, experimenting with controls, asking peers, or waiting to see if the system acts. This interpretive work becomes salient during misalignment, such as when a window blind lowers unexpectedly or a light turns off mid-task. But it also surfaces more subtly: a glance to check for response, or a pause to see if automation will kick in. 

\paragraph{\textbf{Enabling Interaction Quality---Smartness Legibility}}  
Building technologies must enable \textbf{Smartness Legibility}, helping occupants perceive and interpret the building's system behaviour as it unfolds, without needing technical explanation. First, technology should offer brief contextual rationales for actions. A display may state \textit{“Blinds lowering to reduce cooling load”} or \textit{“Window closed due to high outdoor pollen”}. These rationales provide just enough context to make the system's logic legible without overwhelming the user. Systems could also support single-tap prompts like \textit{``Why did this happen?''} or glance-able notifications can help users stay oriented when their assumptions are challenged. 

Importantly, buildings should install \textit{intentional probes}---small, everyday interaction points that allow users to make sense of the building system in action. These may include ``small buttons'' integrated into window frames, gesture-triggered diagnostics (e.g. raising a hand to reveal whether lighting is motion-activated), or ambient indicators such as LEDs, glyphs, or e-ink icons that reveal system activity or invite questions through ``why?” affordances. These are discoverable front-end touchpoints through which occupants can probe the invisible logics driving the environment. Visualisations can further make ongoing or upcoming system actions perceptible. Ambient CO\textsubscript{2} displays~\cite{margariti2024evaluating} in meeting rooms could show rising CO\textsubscript{2} levels as well as indicate when ventilation will increase, which may momentarily introduce drafts or ventilation fan noise. By letting occupants ‘‘ask’’ the building how it works and preview upcoming actions, these probes foster dialogue that helps people anticipate system behaviour and stay in control.


% \begin{tcolorbox}[colback=blue!5!white, colframe=blue!60!black, title=Enabling Interaction Quality: Smartness Legibility]
% Smartness Legibility involves helping occupants perceive and interpret system behaviour as it unfolds, without needing technical explanation or full visibility into system backend logic.
% \end{tcolorbox}

\designbox
{Enabling Interaction Quality: Smartness Legibility}
{Smartness Legibility involves helping occupants perceive and interpret system behaviour as it unfolds, without needing technical explanation or full visibility into system backend logic.}
{Boxed text presenting a conceptual definition}


\subsection{AP6---Settle and Stabilise}
We define \textit{Settle and Stabilise} as the practice of being in a spatial fit---a moment when the space supports one’s needs well enough that active negotiation fades. People stop searching, because the space feels ``good enough.'' Signs of settling include decreased movement, sustained engagement, and fewer glances toward environmental sources of discomfort or disruption. Settling is about reaching a fluency of activity where the space ``gets out of the way.'' Yet this state remains contingent. When conditions shift, such as when a new person enters, the process of adjustment may begin again.

\paragraph{\textbf{Enabling Interaction Quality---Active Non-Intervention}}
To support this practice, technologies must recognise and preserve the tacit equilibrium of settled use. We refer to this as \textbf{Active Non-Intervention}---the system’s capacity to quietly sustain supportive conditions without reasserting control. When signs of settling are detected, such as stillness or minimal user interaction with surrounding automation (e.g. changing the heights on standing desks), systems should defer non-urgent automation (e.g. keeping light levels fixed) to avoid disturbing bodily ease or attentional flow. Systems might anchor the current configuration, offering options like ``save this setup'' until the user appears unsettled. Finally, occasional, affective prompts to mark when ``this just works''---can support slow learning and personalisation without disrupting task immersion. Such cues help both user and system learn what stability feels like, while preserving the quietness that makes it valuable.

% \begin{tcolorbox}[colback=blue!5!white, colframe=blue!60!black, title=Enabling Interaction Quality: Active Non-Intervention]
% Active Non-Intervention is the system’s capacity to recognise and thereby quietly sustain supportive conditions without reasserting control.
% \end{tcolorbox}

\designbox
{Enabling Interaction Quality: Active Non-Intervention}
{Active Non-Intervention is the system’s capacity to recognise and thereby quietly sustain supportive conditions without reasserting control.}
{Boxed text presenting a conceptual definition}


\begin{table*}[htbp]
\centering
\caption{Overview of \textit{Affective Practices} (AP), enabling \textit{Interaction Qualities} (IQ), and corresponding themes. Each row summarises how occupants interact with and interpret smart environments, along with example technologies to support these practices.}
\resizebox{\textwidth}{!}{
\begin{tabular}{p{0.19\textwidth} p{0.23\textwidth} p{0.23\textwidth} p{0.23\textwidth} p{0.08\textwidth}}
\toprule
\textbf{Affective Practice (AP)} & \textbf{Description} & \textbf{Enabling Interaction Quality (IQ)} & \textbf{Example Technologies} & \textbf{Linked Theme(s)} \\
\midrule

\textbf{AP1---Orient \& Appraise} & How users quickly sense whether a space feels ``good enough'' for their current state or task. & \textbf{IQ1---Somatic Legibility}: Make bodily feedback perceptible and interpretable during first impressions. & 
Thermal silhouettes of posture;  
Ambient lighting shifts for CO\textsubscript{2}/heat cues & 
W1 \\

\textbf{AP2---Select a Micro-Condition} & How users select the best nearby space to use or inhabit by comparing small environmental or social differences. & \textbf{IQ2---Spatial Niche Discovery}: Surface contrasts, previews, and enable smooth exploration between spatial pockets. & 
Edge lighting showing noise levels; Live thermal/acoustic gradient maps; Real-time niche ``atmosphere'' tags (e.g. low distraction/15 min) & 
W1, W2, W3  \\

\textbf{AP3---Social Attunement \& Coordination} & How users align their behaviour with evolving social norms and atmospheres. & \textbf{IQ3---Social Cue Surfacing}: Help users sense ongoing social dynamics in shared environments. & 
Sound-reflective light; Corridor activity visualisations; Live ‘atmosphere’ tags (e.g. ``low conversation'') & 
W2, W3  \\

\textbf{AP4---Modulate \& (De)Escalate} & How users regulate discomfort over time through subtle or layered responses. & \textbf{IQ4---(De)Escalation-Aware Interaction}: Offer minimal suggestive support during affective drift. & 
Headphone-based ambient scenes (e.g. “Focus,” “Reset”);  
Scene suggestions triggered by behaviour & 
W1, W4  \\

\textbf{AP5---Probe \& Align with Smart Building Technology} & How users make sense of invisible automation and adjust based on what they learn. & \textbf{IQ5---Smartness Legibility}: Expose system logic through rationales, explanations, \& probeable surfaces. & 
System rationale displays;  
“Why?” buttons or gesture diagnostics;  
Interfaces that invite exploration & 
W5 \\

\textbf{AP6---Settle \& Stabilise} & How users inhabit a good-enough configuration and shift to low-vigilance use. & \textbf{IQ6---Active Non-Intervention}: Recognise settled states and avoid disrupting them. & “Save this setup” checkpoints;  
Automation deferral logic;  
Non-intrusive “this just works” prompts & 
W1, W4 \\
\bottomrule
\end{tabular}
}
\label{tab:ap-iq-summary}
\end{table*}


Our examples span a range of modalities (auditory, visual, spatial), temporal characteristics (continuous vs. discrete and event-triggered), and spatial scales (personal, semi-public, and public). Although we did not begin with a predefined design space, this breadth emerged intuitively from the practices themselves: some practices (e.g. AP4 and AP6) naturally invite personal, bodily, or headphone-based interventions, whereas others (e.g. AP3 and AP5) require public-facing or shared ambient cues to support social attunement or system legibility. Reflecting on our examples reveals that AfI-informed design opportunities necessarily cut across these dimensions, and that affective sense-making in buildings is multi-scalar and multi-modal. We see this as an early mapping rather than an exhaustive typology, and an opportunity for future work to more systematically explore this design space.